\section*{الفصل \ref{ch:b2:structures} --- المجموعات والبنى}

\begin{solution}{exo:b2:structures:1}
\emph{الأجزاء المنتهية من $\N$:} \omterm{def:b2:structures:countable}{قابلة للعد} --- فمجموعة أجزاء
$\intint{0}{n}$ منتهية، والأجزاء المنتهية تشكّل الاتحاد القابل للعد
على $n$ لهذه المجموعات
(\cref{prop:b2:structures:countablestable}~(3))؛ وهي غير منتهية لأنها
تحتوي كل المجموعات الأحادية.

\emph{كل أجزاء $\N$:} غير \omterm{def:b2:structures:countable}{قابلة للعد}، بمبرهنة كانتور
(\cref{thm:b2:structures:cantor}~(1) مع $E = \N$).

\emph{$\R \setminus \Q$:} غير \omterm{def:b2:structures:countable}{قابلة للعد} --- وإلا لكانت $\R = \Q \cup
(\R\setminus\Q)$ اتحاد مجموعتين قابلتين للعد،
وهذا يناقض \cref{thm:b2:structures:cantor}~(2).

\emph{كثيرات الحدود على $\Q$:} \omterm{def:b2:structures:countable}{قابلة للعد} --- فكثيرات الحدود من الدرجة
$\leq n$ تنغرس في $\Q^{n+1}$ (جداءات منتهية لمجموعات \omterm{def:b2:structures:countable}{قابلة للعد})،
ثم نأخذ الاتحاد على $n$.

\emph{المتتاليات الثنائية المنعدمة ابتداءً من رتبة ما:} \omterm{def:b2:structures:countable}{قابلة للعد} ---
فهي تتقابل مع الأجزاء المنتهية من $\N$ (الحامل).
\end{solution}

\begin{solution}{exo:b2:structures:2}
نحسب عنصرًا عنصرًا، مع تطبيق العامل الأيمن أولًا.
يرسل $\sigma\tau$ العناصر $1 \mapsto \sigma(1) = 4$، $\;2 \mapsto \sigma(3)
= 5$، $\;3 \mapsto \sigma(7) = 7$، $\;4 \mapsto \sigma(4) = 2$، $\;5
\mapsto \sigma(5) = 3$، $\;6 \mapsto \sigma(6) = 1$، $\;7 \mapsto
\sigma(2) = 6$:
\[
\sigma\tau = (1\,4\,2\,5\,3\,7\,6),
\]
أي دورة ذات الطول $7$. وكذلك يرسل $\tau\sigma$ العناصر $1 \mapsto \tau(4) = 4$،
$\;2 \mapsto \tau(6) = 6$، $\;3 \mapsto \tau(5) = 5$، $\;4 \mapsto
\tau(2) = 3$، $\;5 \mapsto \tau(3) = 7$، $\;6 \mapsto \tau(1) = 1$،
$\;7 \mapsto \tau(7) = 2$:
\[
\tau\sigma = (1\,4\,3\,5\,7\,2\,6),
\]
أي دورة ذات الطول $7$ أيضًا (كما هو متوقع: $\sigma\tau$ و$\tau\sigma$
متقارنان، ومن ثم لهما نمط الدورات نفسه).

الرتب والإشارات: نمط دورات $\sigma$ هو $(4,2)$: الرتبة
$\operatorname{lcm}(4,2) = 4$، والإشارة $(-1)^3(-1)^1 = +1$؛ و$\tau$
دورة ذات الطول $3$: الرتبة $3$، والإشارة $+1$؛ والجداءان دورتان
ذواتَا الطول $7$: الرتبة $7$، والإشارة $(-1)^6 = +1$.

$\sigma^{2026}$: $2026 = 4 \times 506 + 2$، ومنه $\sigma^{2026} =
\sigma^2 = (1\,2)(4\,6)$ (نربّع \omterm{def:b2:structures:sn}{الدورة} ذات الطول $4$؛ أما \omterm{def:b2:structures:sn}{المبادلة}
فيزيلها التربيع).
\end{solution}

\begin{solution}{exo:b2:structures:3}
$\{0,1\}^{\N} \to \mathcal{P}(\N)$: نرسل كل متتالية إلى حاملها
--- وهذا تقابل (بالدوال المميّزة)، ولا حاجة إلى أي مبرهنة.

$\{0,1\}^{\N} \to \intcc{0}{1}$: التطبيق بالأساس $3$، أي $(a_n) \mapsto
\sum 2a_n 3^{-n-1}$، متباين (فمتتاليتان مختلفتان تختلفان
أول ما تختلفان عند الرتبة $N$؛ ولا تستطيع الذيول تعويض فرق قدره $2\cdot
3^{-N-1}$، لأن $\sum_{n > N} 2\cdot 3^{-n-1} = 3^{-N-1} <
2\cdot3^{-N-1}$).

$\intcc{0}{1} \to \{0,1\}^{\N}$: النشر الثنائي، مع اختيار (مثلًا)
النشر الذي لا ينتهي بالرقم $1$ المتكرر: وهو متباين.

وبمبرهنة كانتور--برنشتاين (\cref{thm:b2:structures:cantorbernstein})
مطبَّقة على الغرستين الأخيرتين، تكون $\intcc{0}{1}$ و
$\{0,1\}^{\N}$ \omterm{def:b2:structures:countable}{متساويتي القوة}، ومنه فالمجموعات الثلاث كلها كذلك.
\end{solution}

\begin{solution}{exo:b2:structures:4}
ليكن $c = ab = ba$ و$d = \operatorname{ord}(c)$. أولًا $c^{mn} =
a^{mn} b^{mn} = e$ (التبادل يسمح بتفكيك القوة)، ومنه $d
\mid mn$. وعكسيًا يعطي $c^d = e$ أن $a^d = b^{-d}$؛ وهذا العنصر
يقع في $\langle a\rangle \cap \langle b\rangle$، وهي زمرة جزئية تقسم
رتبتها كلًا من $m$ و$n$ (لاغرانج في كل \omterm{def:b2:structures:generated}{زمرة دورية})،
ومن ثم فهي محايدة: $a^d = b^d = e$، ومنه $m \mid d$ و$n \mid d$،
وبالأولية فيما بينهما $mn \mid d$. ومنه $d = mn$.

وفي $\mathfrak{S}_3$: نأخذ $a = (1\,2)$ (رتبتها $2$) و$b =
(1\,2\,3)$ (رتبتها $3$)، ورتبتاهما أوليتان فيما بينهما، وهما لا تتبادلان: فرتبة $ab =
(2\,3)$ هي $2 \neq 6$ --- وفعلًا ليس في $\mathfrak{S}_3$
عنصر رتبته $6$. فالتبادل شرط جوهري.
\end{solution}

\begin{solution}{exo:b2:structures:5}
نقابل بين كل $x \in G$ و$x^{-1}$. الأزواج $\{x, x^{-1}\}$ التي
$x \neq x^{-1}$ ذات عنصرين وتقسم اتحادها؛ والعناصر
الباقية هي بالضبط تلك التي $x = x^{-1}$، أي $x^2 =
e$. وبما أن $\abs G$ زوجي وأن الأزواج ذات العنصرين تغطي عددًا
زوجيًا من العناصر، فإن قوة المجموعة $\{x : x^2 = e\}$ زوجية؛
وهي تحتوي $e$، ومن ثم تحتوي عنصرًا آخر $x \neq
e$ على الأقل --- أي عنصرًا رتبته $2$.
\end{solution}

\begin{solution}{exo:b2:structures:6}
كل عنصر من $A_n$ جداء عدد زوجي من
المبادلات (\cref{thm:b2:structures:signature}: نفكك إلى
مبادلات؛ والعدد زوجي لأن الإشارة $+1$). ويكفي
أن نكتب كل جداء مبادلتين بدورات ذات الطول $3$:
\[
(a\,b)(a\,c) = (a\,c\,b),
\qquad
(a\,b)(c\,d) = (a\,c\,b)(a\,c\,d) \quad (\text{مختلفة } a,b,c,d),
\]
(بالتحقق عبر التقييم)، و$(a\,b)(a\,b) = \mathrm{id}$. ومنه فالدورات
ذات الطول $3$ تولّد $A_n$.
\end{solution}

\begin{solution}{exo:b2:structures:7}
\emph{$(\Q,+) \to (\Z,+)$:} التشاكل المعدوم وحده. فمن أجل أي $x$ ومن أجل
كل $n \geq 1$، يقبل $f(x) = n f\bigl(\frac xn\bigr)$ القسمة
على $n$ في $\Z$؛ والعدد الصحيح الوحيد الذي يقبل القسمة على كل $n$ هو $0$، ومنه
$f(x) = 0$ لكل $x$.

\emph{$(\Z/n\Z, +) \to (\Z/m\Z, +)$:} يتحدد التشاكل بالقيمة
$c = f(\overline 1)$، التي يجب أن تحقق $n c \equiv 0 \pmod m$،
أي أن $c$ مضاعف للعدد $\frac{m}{\gcd(m,n)}$؛ وهناك
$\gcd(m,n)$ صنفًا كهذه، وكل اختيار يعرّف تشاكلًا فعلًا
(نحلّل $k \mapsto kc$ عبر $\Z/n\Z$ بالخاصية الشمولية).

\emph{$(\Q, +) \to (\Q_+^*, \times)$:} التشاكل المحايد وحده. فإذا كان
$f(x) = y$، فمن أجل كل $n$ لدينا $y = f(n \cdot \frac xn) =
f(\frac xn)^n$، أي إن هذه القيمة قوة $n$-ية في $\Q_+^*$. لكن عددًا ناطقًا
$y \neq 1$ لا يمكن أن يكون قوة $n$ من أجل كل $n$: إذ يظهر عدد أولي ما
في $y$ بأس غير معدوم $v$، و$n \nmid v$ من أجل $n > \abs
v$ (فأسس القوى ذوات الأس $n$ مضاعفات للعدد $n$، بحكم وحدانية
التفكيك). ومنه $f \equiv 1$.
\end{solution}

\begin{solution}{exo:b2:structures:8}
$360 = 2^3 \cdot 3^2 \cdot 5$:
$\varphi(360) = 360\bigl(1 - \tfrac12\bigr)\bigl(1 -
\tfrac13\bigr)\bigl(1 - \tfrac15\bigr) = 360 \cdot \tfrac12 \cdot
\tfrac23 \cdot \tfrac45 = 96$.

الجملة: القياسات $8, 9, 5$ أولية فيما بينها مثنى مثنى، وجداؤها $360$. من $x
\equiv 3 \pmod 8$ و$x \equiv 5 \pmod 9$: نجد $x = 3 + 8k$ مع $3 +
8k \equiv 5 \pmod 9$، أي $-k \equiv 2$، و$k \equiv -2 \equiv 7
\pmod 9$: ومنه $x \equiv 3 + 56 = 59 \pmod{72}$. ثم $59 + 72\ell \equiv
2 \pmod 5$: $4 + 2\ell \equiv 2$، $2\ell \equiv 3 \equiv 8$، $\ell
\equiv 4 \pmod 5$: ومنه $x \equiv 59 + 288 = 347 \pmod{360}$.

الرقمان الأخيران من $3^{2026}$: بترديد $4$، $3^{2026} = 9^{1013} \equiv
1$. وبترديد $25$: $\varphi(25) = 20$ و$2026 = 20\cdot101 + 6$، ومنه
$3^{2026} \equiv 3^6 = 729 \equiv 4 \pmod{25}$. نحلّ $x \equiv 1
\pmod 4$، $x \equiv 4 \pmod{25}$: يعطي $x = 4 + 25k \equiv 1 \pmod 4$
أن $k \equiv 1 \pmod 4$: ومنه $x \equiv 29 \pmod{100}$. والرقمان الأخيران
هما $29$.
\end{solution}

\begin{solution}{exo:b2:structures:9}
لتكن $A$ حلقة تامة منتهية وليكن $a \in A$، $a \neq 0$.
التطبيق $x \mapsto ax$ متباين ($ax = ay \implies a(x - y) = 0
\implies x = y$، إذ لا قواسم للصفر)؛ وكل تطبيق متباين من مجموعة منتهية
إلى نفسها يكون شاملًا (مجلد السنة الأولى، تكافؤ مبدأ الجحور).
ومنه $1 = ab$ من أجل $b$ ما: فكل عنصر غير معدوم قابل للقلب، و$A$
حقل.

$\Z/n\Z$: إذا كان $n$ أوليًا فهي حلقة تامة ($n \mid ab
\implies n \mid a$ أو $n \mid b$، بمبرهنة إقليدس المساعدة)، ومنتهية، ومن ثم
حقل؛ وإذا كان $n = rs$ مركّبًا فإن $\overline r\,\overline s =
\overline 0$ يُبرز قواسم للصفر.
\end{solution}

\begin{solution}{exo:b2:structures:10}
ليكن $m = \max\{\operatorname{ord}(x) : x \in G\}$، وهي مبلوغة عند $a$.

\emph{الادعاء: رتبة كل $x \in G$ تقسم $m$.} لنفترض أن رتبة
$x$ ما هي $q$ مع $q \nmid m$: عندئذٍ توجد قوة أولية $p^k$
تقسم $q$ ولا تقسم $m$. لنكتب $m = p^j m'$ مع $p \nmid m'$ و$j
< k$. رتبة العنصر $a^{p^j}$ هي $m'$؛ ورتبة العنصر
$x^{q/p^k}$ هي $p^k$؛ والرتبتان أوليتان فيما بينهما والعنصران
يتبادلان (لأن $G \subseteq K^*$ تبادلية)، ومنه حسب
\cref{exo:b2:structures:4} تكون رتبة جدائهما $p^k m' > p^j m'
= m$: وهذا يناقض الأعظمية.

فتحقق إذن كل $x \in G$ العلاقة $x^m = 1$: أي أن لكثير الحدود $X^m - 1$
$\abs G$ جذرًا على الأقل في الحقل $K$، ومنه $\abs G \leq m$
(فكثير حدود غير معدوم من الدرجة $m$ له $m$ جذرًا على الأكثر، مجلد
السنة الأولى). لكن $m = \operatorname{ord}(a) \leq \abs G$ بمبرهنة لاغرانج.
ومنه $m = \abs G$، وتكون $\langle a \rangle$، وقوتها $m =
\abs G$، هي $G$ كلها: أي إنها دورية.

ومن أجل $K = \Z/p\Z$: تكون $(\Z/p\Z)^*$ زمرة جزئية منتهية من $K^*$، ومن ثم
دورية (رتبتها $p - 1$).
\end{solution}

\begin{solution}{exo:b2:structures:11}
\emph{ليست دورية:} تتكوّن الزمرة الجزئية $\langle \frac pq\rangle$
من المضاعفات الصحيحة للعدد $\frac pq$، وكلها ذات
مقام يقسم $q$ (بعد الاختزال)؛ ومن ثم لا تحتوي
$\frac{1}{2q}$. فلا يستطيع مولّد وحيد أن يبلغ المقامات غير
المحدودة في $\Q$.

\emph{ليست مولَّدة بجزء منته:} تتكوّن الزمرة الجزئية \omterm{def:b2:structures:generated}{المولَّدة}
بالعناصر $\frac{p_1}{q_1}, \dots, \frac{p_k}{q_k}$ من الأعداد الناطقة
التي تقسم مقاماتها $Q = q_1 \cdots q_k$ (فالتراكيب الصحيحة
لها مقام يقسم $Q$): وهي لا تحتوي $\frac{1}{2Q}$.

\emph{الزمر الجزئية \omterm{def:b2:structures:generated}{المولَّدة} بجزء منته دورية:} مع $Q$ كما تقدم،
تكون الزمرة الجزئية $H = \langle \frac{p_1}{q_1}, \dots,
\frac{p_k}{q_k}\rangle$ محتواة في $\frac{1}{Q}\Z$. والتطبيق $x
\mapsto Qx$ تماثل من $\frac1Q\Z$ على $\Z$ يحمل
$H$ إلى زمرة جزئية من $\Z$، وهي $n\Z$ من أجل $n$ ما (مجلد
السنة الأولى): ومنه فالزمرة $H = \frac{n}{Q}\Z$ دورية، مولَّدة بالعنصر
$\frac nQ$.
\end{solution}

\begin{solution}{exo:b2:structures:12}
\emph{جزء قابل للعد.} لتكن $E$ غير منتهية. نبني $a_0,
a_1, a_2, \dots$ بالتراجع: $E$ غير خالية، فنختار $a_0 \in E$؛
وإذا اختيرت $a_0, \dots, a_n$، فإن $E \setminus \{a_0, \dots,
a_n\}$ غير خالية (لأن $E$ غير منتهية)، فنختار $a_{n+1}$ فيها. والعناصر
$a_n$ مختلفة مثنى مثنى بحكم الإنشاء، ومنه فإن $A = \{a_n : n
\in \N\}$ جزء قابل للعد من $E$.

\emph{مجموعة غير منتهية $\implies$ مساوية القوة لجزء صحيح منها.} نعرّف
$f \colon E \to E \setminus \{a_0\}$ بالعلاقتين $f(a_n) = a_{n+1}$ و
$f(x) = x$ من أجل $x \notin A$. وهو متباين (فالقطعتان
متباينتان وصورتاهما منفصلتان) وشامل على $E \setminus
\{a_0\}$: فكل $a_{n+1}$ مبلوغ، وكل $x \notin A$ مبلوغ. ومنه فإن
$E$ مساوية القوة للجزء الصحيح $E \setminus \{a_0\}$.

\emph{العكس.} إذا كانت $E$ منتهية وكان $g \colon E \to F$ تقابلًا
على $F \subseteq E$ مع $F \neq E$، فإن $g$ غرسة
من $E$ في نفسها وليست شاملة،
وهذا يناقض مبدأ الجحور (مجلد السنة الأولى: كل تطبيق ذاتي
متباين على مجموعة منتهية تقابلي). ومنه فكل مجموعة
مساوية القوة لجزء صحيح منها تكون غير منتهية.
\end{solution}

\begin{solution}{pb:b2:structures:1}
\textbf{1.} تُسند الوضعية إلى كل خانة من الخانات $16$ محتوى
من المحتويات $16$ (القطع $1$--$15$ أو الفراغ $b = 16$)،
مرة واحدة بالضبط لكل منها: أي بالضبط تقابل $\intint1{16} \to
\intint1{16}$، وهو عنصر من $\mathfrak{S}_{16}$؛ وعددها $16!
= 20\,922\,789\,888\,000$. وتُزلق النقلة المشروعة قطعة واحدة
مجاورة للفراغ، فيكون عدد النقلات هو عدد جيران خانة
الفراغ: $2$ من أجل الخانات الزاوية الأربع،
و$3$ من أجل الخانات الحافية الثماني، و$4$ من أجل الخانات الداخلية الأربع.

\textbf{2.} بعد الانزلاق، تحمل الخانة $p$ المحتوى السابق للخانة
$c$ وتحمل الخانة $c$ الفراغ؛ وتبقى كل الخانات الأخرى دون تغيير:
$\sigma'(p) = \sigma(c)$، $\sigma'(c) = \sigma(p) = 16$،
و$\sigma' = \sigma$ في ما عدا ذلك. وهذا بالضبط $\sigma' = \sigma
\circ (p\ c)$. وبما أن $\varepsilon$ تشاكل وأن
$\varepsilon\bigl((p\ c)\bigr) = -1$:
$\varepsilon(\sigma') = -\varepsilon(\sigma)$.

\textbf{3.} تختلف الخانتان المتجاورتان بخطوة واحدة في إحدى
الإحداثيتين دون الأخرى، ومن ثم يتغير تماثل $i + j$: فتأخذ $\chi$
قيمتين متعاكستين على خانتين متجاورتين. وتنقل النقلةُ الفراغ
من $p$ إلى $c$ المجاورة، فتقلب $\chi(\text{خانة الفراغ})$.
وعلى امتداد مسير مغلق للفراغ، تنقلب $\chi$ مرة واحدة في كل نقلة
وتعود إلى قيمتها الابتدائية: فعدد النقلات زوجي.

\textbf{4.} حسب السؤالين 2 و3، تقلب النقلة الواحدة عاملَي
$I(\sigma) = \varepsilon(\sigma)\chi(\sigma^{-1}(16))$ كليهما؛ فلا يتغير
جداؤهما. ومن أجل الوضعية المحلولة:
$\varepsilon(\mathrm{id}) = +1$ والفراغ في الخانة $16$،
أي السطر $4$ والعمود $4$: ومنه $\chi(16) = (-1)^{8} = +1$، فيكون
$I(\mathrm{id}) = +1$.

\textbf{5.} $\sigma_L$ هي \omterm{def:b2:structures:sn}{مبادلة} الخانتين $(14\ 15)$:
$\varepsilon(\sigma_L) = -1$؛ وفراغها في موطنه، $\chi(16) = +1$:
ومنه $I(\sigma_L) = -1 \neq +1 = I(\mathrm{id})$. وبما أن $I$
تصمد أمام كل نقلة، فلا متتالية نقلات تربط $\sigma_L$
بالوضعية $\mathrm{id}$. فالجائزة كانت في أمان بنيوي.

\textbf{6.} نثبّت خانة $p$ وخانتين أخريين $c \neq d$
مختلفتين عن $p$، ونضع $\tau_0 = (c\ d)$. وعلى مجموعة
الوضعيات ذات الفراغ في $p$، يكون التطبيق $\sigma \mapsto \sigma
\circ \tau_0$ تقابلًا ذاتيًا من الرتبة اثنين (فهو يحافظ على $\sigma(p) = 16$
لأن $\tau_0$ يثبّت $p$) ويقلب $\varepsilon$، ومن ثم يقلب
$I$: فهو يقابل الوضعيات ذات $I = +1$ تقابلًا مع
تلك ذات $I = -1$. ومنه فكل موضع من مواضع الفراغ $16$
يسهم بعدد $15!/2$ من الوضعيات تحقق $I = +1$، ويكون
\[
\abs{\{I = +1\}} = 16 \cdot \frac{15!}{2} = \frac{16!}{2}.
\]

\textbf{7.} النقلة التي تُزلق قطعة $c$ إلى $p$ تُلغى
بإزلاق القطعة نفسها (وهي الآن في $p$) عائدةً إلى $c$: فالتركيب
مع $(p\ c)$ مرتين هو المطابق. ومنه: الانعكاسية (بالمتتالية
الخالية)، والتناظر (بعكس المتتالية وإلغاء كل نقلة)،
والتعدي (بالتلاصق): أي علاقة تكافؤ. وكل
$\sigma \in R$ يحقق $I(\sigma) = I(\mathrm{id}) = +1$ حسب السؤال
4، ومنه $R \subseteq \{I = +1\}$؛ و$\sigma_L \notin R$ يعطي
صنفًا ثانيًا.

\textbf{8.} إذا كان $\sigma(16) = 16$، فإن $\sigma$ يبدّل
الخانات $1, \dots, 15$؛ ولنُسمّ $\rho$ هذا المقصور. وإضافة
نقطة ثابتة لا تغيّر نمط الدورات ولا الإشارة
(نفكك $\rho$ إلى مبادلات؛ والجداء نفسه صالح في
$\mathfrak{S}_{16}$)، ومنه $\varepsilon(\sigma) =
\varepsilon(\rho)$، ويعطي $\chi(16) = +1$ أن $I(\sigma) =
\varepsilon(\rho)$. وكل وضعية يمكن نقلها إلى وضعية
فراغها في موطنه: فالشبكة مترابطة، فنسير بالفراغ على
مسير من خانات متجاورة إلى الخانة $16$ (وكل خطوة نقلة مشروعة).
ولنفترض الآن أن كل $\rho \in \mathfrak{S}_{15}$ زوجية محقَّقة
ببرنامج. فلتكن $\sigma$ مع $I(\sigma) = +1$: نسير بالفراغ
إلى موطنه فنبلغ $\widetilde\sigma$ (وهي مكافئة للوضعية $\sigma$)،
مع $I(\widetilde\sigma) = +1$، أي إن مقصورها $\rho$
زوجي؛ والبرنامج الذي يحقق $\rho$ ينقل $\widetilde\sigma$ إلى
$\widetilde\sigma \circ \rho^{-1} = \mathrm{id}$ (انظر السؤال
9). وبالتعدي $\sigma \in R$، ومنه $\{I = +1\} \subseteq
R$ والمساواة.

\textbf{9.} \emph{نقلة واحدة:} ينتهي محتوى $c$ في $p$
والفراغ في $c$: فالأثر هو $\pi = (p\ c)$، وفعلًا
$\sigma' = \sigma \circ (p\ c) = \sigma \circ \pi^{-1}$.
\emph{بالتراجع:} إذا كان لمتتالية أثر $\pi_1$ وكانت تنقل
$\sigma$ إلى $\sigma \circ \pi_1^{-1}$، فإن إتباعها بنقلة
أثرها $\pi_2 = (p'\ c')$ يعطي $(\sigma \circ \pi_1^{-1})
\circ \pi_2^{-1} = \sigma \circ (\pi_2\pi_1)^{-1}$، وتنتقل المحتويات
بالتبديلة $\pi_2 \circ \pi_1$ (أولًا $\pi_1$، ثم $\pi_2$). فالآثار
تتركّب، والبرنامج المنفَّذ انطلاقًا من $\sigma$ ينتهي عند $\sigma
\circ \pi^{-1}$. \emph{الزمرة الجزئية:} أثر البرنامج الخالي هو
$\mathrm{id}$؛ والتلاصق يعطي الجداءات؛ وعكس البرنامج
(السؤال 7) يعطي المقلوبات. وأثر البرنامج يثبّت الخانة $16$
(فالفراغ يبدأ وينتهي في موطنه)، ومنه $H \leq \mathfrak{S}_{15}$.
\emph{الزوجية:} برنامج من $k$ نقلة له $k$ زوجي (السؤال
3)، ويفرض $\varepsilon(\sigma \circ \pi^{-1}) =
(-1)^k\varepsilon(\sigma)$ أن $\varepsilon(\pi) = +1$: ومنه $H
\subseteq A_{15}$.

\textbf{10.} نتتبع الانزلاقات الأربعة انطلاقًا من الفراغ في $16$: النقلة
$16 \to 12$ ترسل محتوى $12$ إلى $16$؛ والنقلة $12 \to 11$
ترسل محتوى $11$ إلى $12$؛ والنقلة $11 \to 15$ ترسل
محتوى $15$ إلى $11$؛ والنقلة $15 \to 16$ ترسل المحتوى
المركون في $16$ (وأصله في $12$) إلى $15$. والحصيلة: $11 \mapsto
12$، $12 \mapsto 15$، $15 \mapsto 11$، والفراغ في موطنه: فالأثر هو
$(11\ 12\ 15)$. والجولة العكسية تلغيه: أثرها $(11\ 12\
15)^{-1} = (11\ 15\ 12)$. وكلاهما أثر برنامج، ومن ثم في
$H$.

\textbf{11.} تجاور الخانات المتتالية: ففي كل زوج مذكور تختلف
الخانتان بمقدار $1$ في السطر نفسه ($16{-}15$،
$15{-}14$، $14{-}13$؛ $1{-}2$، $2{-}3$، $3{-}4$؛ $8{-}7$،
$7{-}6$؛ $10{-}11$، $11{-}12$) أو بمقدار $4$ داخل عمود
($13{-}9$، $9{-}5$، $5{-}1$؛ $4{-}8$؛ $6{-}10$؛ $12{-}16$): أي
مسير مغلق عبر الخانات $16$ كلها، طوله $16$. أما الأثر: فكما
في السؤال 10، بكتابة الخانات المزارة $c_0 = 16, c_1 = 15,
\dots, c_{15} = 12$: ينتقل محتوى $c_i$ إلى $c_{i-1}$ من أجل
$i = 2, \dots, 15$، ويُحمل محتوى $c_1$، المركون في $16$
بعد النقلة الأولى، إلى $c_{15}$ بالنقلة الأخيرة.
فيرسل الأثر إذن $15 \mapsto 12$، و$14 \mapsto 15$، $13
\mapsto 14$، $9 \mapsto 13$، $5 \mapsto 9$، $1 \mapsto 5$، $2
\mapsto 1$، $3 \mapsto 2$، $4 \mapsto 3$، $8 \mapsto 4$، $7
\mapsto 8$، $6 \mapsto 7$، $10 \mapsto 6$، $11 \mapsto 10$، $12
\mapsto 11$: أي بالضبط \omterm{def:b2:structures:sn}{الدورة} ذات الطول $15$، وهي $\zeta$. ويبدأ ترتيبها الدوري
بالخانة $x_0 = 15$، $x_1 = 12$، $x_2 = 11$، ويرسل $(x_0\ x_1\ x_2)
= (15\ 12\ 11)$ العنصر $15 \mapsto 12 \mapsto 11 \mapsto 15$ ---
وهو بالضبط $(11\ 15\ 12)$، أي الجولة الأولية العكسية.

\textbf{12.} ليكن $\gamma = (a\ b\ c)$ و$x \in \intint1n$. إذا كان
$x = g(a)$: فإن $g\gamma g^{-1}(x) = g(\gamma(a)) = g(b)$؛ وكذلك
$g(b) \mapsto g(c)$ و$g(c) \mapsto g(a)$. وإذا كان $x \notin
\{g(a), g(b), g(c)\}$، فإن $g^{-1}(x) \notin \{a,b,c\}$
مثبَّت بالتبديلة $\gamma$، ومن ثم يبقى $x$ ثابتًا. ومنه $g\gamma g^{-1} =
(g(a)\ g(b)\ g(c))$. ومن أجل $g, h \in H$، نجد $ghg^{-1} \in H$ بحكم
بديهيات الزمرة الجزئية.

\textbf{13.} $\zeta \in H$ (السؤال 11) و$s_0 = (x_0\ x_1\
x_2) \in H$ (السؤالان 10--11). وبما أن $\zeta(x_i) = x_{i+1}$
(بأدلة بترديد $15$)، يعطي السؤال 12 أن
\[
\zeta^{t}\,s_0\,\zeta^{-t}
= \bigl(\zeta^t(x_0)\ \zeta^t(x_1)\ \zeta^t(x_2)\bigr)
= (x_t\ x_{t+1}\ x_{t+2}) = s_t \in H
\qquad (t = 0, 1, \dots, 14).
\]

\textbf{14.} إلى حدّ القلب، لنفترض أن $s = (a\ b\ c)$ و$t =
(b\ c\ d)$ (\omterm{def:b2:structures:sn}{فالدورة} ذات الطول $3$ على $\{a,b,c\}$ هي $(a\ b\ c)$ أو
مقلوبها؛ وكذلك على $\{b,c,d\}$؛ وتعويض مولّد بمقلوبه
لا يغيّر $\langle s, t\rangle$). عندئذٍ، بتطبيق
$t$ أولًا،
\[
st \colon a \mapsto b,\quad b \mapsto a,\quad c \mapsto d,\quad
d \mapsto c, \qquad\text{أي}\quad st = (a\ b)(c\ d),
\]
وهي \omterm{def:b2:structures:sn}{مبادلة} مضاعفة. وتتكوّن الزمرة الجزئية $G = \langle s, t\rangle$
من التبديلات الزوجية للحروف الأربعة، ومنه $G \leq
A_4$ و$\abs G \mid 12$؛ وهي تحتوي عنصرًا رتبته $3$
وآخر رتبته $2$، ومن ثم $6 \mid \abs G$ (لاغرانج،
\cref{thm:b2:structures:lagrange}، مطبَّقة على الزمرتين الدوريتين
الجزئيتين). ولو كانت في $A_4$ زمرة جزئية $K$ رتبتها $6$، لكان
دليلها $2$، ولكان عندئذٍ $g^2 \in K$ من أجل كل $g \in A_4$: فمن أجل
$g \in K$ هذا واضح؛ ومن أجل $g \notin K$ لا يوجد إلا الصنفان الجانبيان
$K$ و$gK$، ومن ثم فالصنف $g^2K$ هو $K$ أو $gK$، ويفرض $g^2K =
gK$ أن $g \in K$. فكل مربّع يقع إذن في $K$. لكن كل
دورة ذات الطول $3$، أي $\gamma$، مربّع، إذ $\gamma = (\gamma^2)^2$،
و$A_4$ تحتوي ثماني دورات ذات الطول $3$: ومنه $8 > 6$، وهو تناقض. ومنه
$\abs G = 12$: أي $G = A_4$.

\textbf{15.} نمدّد $u \mapsto a$ و$v \mapsto b$ إلى تقابل
$g_0$ على $X$ (بإرسال الحروف $k - 2$ الباقية تقابليًا
إلى متممة $\{a, b\}$ كيفما اتفق). وإذا كانت $g_0$ فردية،
نختار حرفين مختلفين $s_1, t_1 \in X \setminus \{u, v\}$
(وهذا ممكن لأن $k \geq 4$) ونعوّض $g_0$ بالتبديلة $g_0 \circ (s_1\
t_1)$، وهي زوجية وما تزال ترسل $u \mapsto a$، $v \mapsto
b$. ثم نمدّد بالمطابق خارج $X$: فنحصل على تبديلة زوجية $g \in
G$ (فهي تبديلة زوجية للمجموعة $X$). عندئذٍ يعطي السؤال 12:
\[
g\,(u\ v\ w)\,g^{-1} = (g(u)\ g(v)\ g(w)) = (a\ b\ w) \in G,
\]
باستعمال $g(w) = w$.

\textbf{16.} كل دورة ذات الطول $3$ في $X \cup \{w\}$ تقع في $G$:
فالدورات المحمولة في $X$ تبديلات زوجية للمجموعة $X$؛ وأما دورة حاملها
$\{a, b, w\}$ فهي $(a\ b\ w)$ أو $(b\ a\ w)$، وكلتاهما
يعطيها السؤال 15. وحسب \cref{exo:b2:structures:6}، تولّد
الدورات ذات الطول $3$ في المجموعة $X \cup \{w\}$ ذات $(k+1)$ عنصرًا
زمرتَها المتناوبة، ومنه فإن $G$ تحتوي كل تبديلة زوجية
للمجموعة $X \cup \{w\}$. \emph{التسلسل:} ليكن $G = \langle s_0, \dots,
s_{12}\rangle$. تعطي المبرهنة المساعدة أ مطبَّقة على $s_0 = (x_0\ x_1\ x_2)$ و
$s_1 = (x_1\ x_2\ x_3)$ (وحاملاهما يشتركان في $\{x_1, x_2\}$)
كل التبديلات الزوجية للمجموعة $X_4 = \{x_0, x_1, x_2, x_3\}$. وإذا كانت $G$
تحتوي كل التبديلات الزوجية للمجموعة $X_m = \{x_0, \dots,
x_{m-1}\}$ ($4 \leq m \leq 14$)، فإن $s_{m-2} = (x_{m-2}\
x_{m-1}\ x_m)$ فيها $u = x_{m-2}, v = x_{m-1} \in X_m$ والحرف
الجديد $w = x_m$: فتعطي المبرهنة المساعدة ب مع الجزء الأول كل التبديلات
الزوجية للمجموعة $X_{m+1}$. وبالتراجع إلى غاية $m = 14$: $G \supseteq
A_{15}$ (كل التبديلات الزوجية للخانات الخمس عشرة)، و$G
\subseteq A_{15}$ لأن كل $s_t$ زوجية: ومنه $\langle s_0, \dots,
s_{12}\rangle = A_{15}$.

\textbf{17.} السؤالان 13 و16: $A_{15} = \langle s_0, \dots,
s_{12}\rangle \subseteq H$؛ والسؤال 9: $H \subseteq A_{15}$. ومنه
$H = A_{15}$، ورتبتها $15!/2 = 653\,837\,184\,000$: فكل إعادة ترتيب
زوجية للقطع الخمس عشرة أثرُ برنامج.

\textbf{18.} أرجع السؤال 8 إثبات $R = \{I = +1\}$ إلى تحقيق
كل $\rho \in \mathfrak{S}_{15}$ زوجية ببرنامج: وقد تمّ ذلك في
السؤال 17. ومع السؤال 6، نجد $\abs R = 16!/2 =
10\,461\,394\,944\,000$. \emph{الصنفان:} ليؤثّر $t_0 = (14\
15)$ على \emph{المحتويات}: $\varphi(\sigma) = t_0 \circ
\sigma$. فالنقلة المشروعة انطلاقًا من $\sigma$ نقلة مشروعة انطلاقًا من
$\varphi(\sigma)$ (فخانة الفراغ لا تتغير:
$(t_0\sigma)^{-1}(16) = \sigma^{-1}(t_0(16)) = \sigma^{-1}(16)$،
والخانة المنقولة هي نفسها)، و$\varphi(\sigma \circ \tau)
= \varphi(\sigma) \circ \tau$: ومنه فإن $\varphi$ يرسل متتاليات النقلات إلى
متتاليات نقلات، تقابليًا (فهو تقابل ذاتي من الرتبة اثنين). وهو يقلب $I$:
$\varepsilon(t_0\sigma) = -\varepsilon(\sigma)$، مع خانة
الفراغ نفسها. ومنه فإن $\varphi$ يرسل صنف $R = \{I = +1\}$
الوضعية $\mathrm{id}$ تقابليًا على صنف $\varphi(\mathrm{id})
= \sigma_L$، وهو من ثم كل $\{I = -1\}$: أي صنفان
اثنان بالضبط. وهذه هي مبرهنة جونسون--ستوري.

\textbf{19.} نرقّم الخانات بترتيب القراءة ولتكن $k = 4(i -
1) + j$ خانة الفراغ. نعدّ انقلابات $\sigma$
(أي أزواج الخانات $x < y$ التي $\sigma(x) > \sigma(y)$): تسهم أزواج
خانتَي قطعتين بمقدار $N$؛ وأما الأزواج التي تشمل الفراغ: فالخانات
التي بعد الفراغ تحمل كلها قطعًا $< 16$، وكل منها منقلب ($16 - k$
زوجًا)، والخانات التي قبله لا تنقلب أبدًا. ومنه
$\varepsilon(\sigma) = (-1)^{N + 16 - k} = (-1)^{N + k}$. وبما أن
$k = 4(i-1) + j \equiv j \pmod 2$،
\[
I(\sigma) = (-1)^{N + j}\,(-1)^{i + j} = (-1)^{N + i}
= (-1)^{N + r + 1}
\]
باستعمال $i = 5 - r$. وحسب السؤال 18، تكون $\sigma$ قابلة للحلّ إذا وفقط إذا
$I(\sigma) = +1$ إذا وفقط إذا كان $N + r$ فرديًا. وللتحقق: الوضعية المحلولة، $N = 0$، $r
= 1$: فردي، فهي قابلة للحلّ؛ ووضعية لويد، $N = 1$، $r = 1$: زوجي، فهي غير قابلة للحلّ.

\textbf{20.} \emph{الفعل:} $e \cdot \sigma = \sigma \circ
\mathrm{id} = \sigma$ و$g \cdot (h \cdot \sigma) = \sigma
\circ h^{-1} \circ g^{-1} = \sigma \circ (gh)^{-1} = (gh) \cdot
\sigma$؛ و$\sigma \circ h^{-1}$ وضعية فراغها في موطنه
من جديد (لأن $h$ يثبّت الخانة $16$). \emph{الحرية:} يعطي $\sigma \circ
h^{-1} = \sigma$ أن $h^{-1} = \mathrm{id}$ (بالتركيب مع
$\sigma^{-1}$). \emph{المدارات = أصناف البرامج:} يقول السؤال 9
إن الوضعيات القابلة للبلوغ من $\sigma$ بالبرامج هي
بالضبط $\sigma \circ \pi^{-1}$، $\pi \in H$: أي المدار $H
\cdot \sigma$. \emph{العدّ:} تجعل الحرية $h \mapsto h \cdot
\sigma$ متباينًا، فيكون لكل مدار $\abs H = 15!/2$ عنصرًا؛
ومن ثم تنقسم الوضعيات $15!$ ذات الفراغ في الموطن إلى
$15!\,/\,(15!/2) = 2$ مدارًا --- وهي ظلّ صنفَي جونسون--ستوري
في الوضعيات ذات الفراغ في الموطن.

\textbf{21.} الشبكة $3 \times 3$ ثنائية التجزئة بتلوين
رقعة الشطرنج: فكل خطوة في مسير تغيّر اللون، ومن ثم
يكون طول كل مسير \emph{مغلق} زوجيًا. ولو وُجد مسير مغلق يزور
كلًا من الخانات $9$ مرة واحدة بالضبط لكان طوله $9$، وهو فردي:
وهذا مستحيل. فإنشاء الجولة الكبرى في الجزء الثالث غير
متاح إذن في أحجية الثمانية.

\textbf{22.} \emph{جولة المحيط} $9 \to 8 \to 7 \to 4 \to 1
\to 2 \to 3 \to 6 \to 9$ (وكل خطواتها بين متجاورين؛ وطولها $8$، وهو زوجي):
حسب حساب السؤال 11 مع $c_1 = 8, c_2 = 7, c_3 =
4, c_4 = 1, c_5 = 2, c_6 = 3, c_7 = 6$، يكون الأثر
\[
\zeta' = (8\ 6\ 3\ 2\ 1\ 4\ 7),
\]
أي دورة ذات الطول $7$ تثبّت المركز $5$ (فمحتوى $7$ ينتقل إلى $8$،
ومحتوى $4$ إلى $7$، ومحتوى $1$ إلى $4$، ومحتوى $2$ إلى $1$، ومحتوى $3$ إلى $2$، ومحتوى
$6$ إلى $3$، ومحتوى $8$ إلى $6$). \emph{جولة الزوايا} $9 \to 6 \to
5 \to 8 \to 9$: أثرها $(6\ 8\ 5)$ (فمحتوى $5$ ينتقل إلى
$6$، ومحتوى $8$ إلى $5$، ومحتوى $6$ --- المركون في $9$ --- إلى $8$). ونضع
$y_t = \zeta'^{\,t}(8)$: فنجد $y_0 = 8, y_1 = 6, y_2 = 3, y_3 = 2,
y_4 = 1, y_5 = 4, y_6 = 7$. والمقارنة (السؤال 12):
\[
\zeta'^{\,t}\,(6\ 8\ 5)\,\zeta'^{-t}
= (y_{t+1}\ y_t\ 5) =: T_t \in H_{3\times3},
\]
لأن $\zeta'$ يثبّت $5$. وحاملا $T_0 = (y_1\ y_0\ 5)$
و$T_1 = (y_2\ y_1\ 5)$ يشتركان بالضبط في $\{y_1, 5\}$: فتعطي المبرهنة المساعدة أ
كل التبديلات الزوجية للمجموعة $\{y_0, y_1, y_2, 5\}$. ثم يضيف
$T_2 = (y_3\ y_2\ 5)$ العنصر $y_3$ بالمبرهنة المساعدة ب (فحروفه
$y_2, 5$ تقع في المجموعة الجارية، و$k = 4$)، وتضيف $T_3, T_4, T_5$
بدورها $y_4, y_5, y_6$: فتقع كل التبديلات الزوجية للخانات
الثماني غير الموطن في زمرة البرامج، التي تتكوّن
هي أيضًا من تبديلات زوجية (فحجة السؤال 9 لا تتعلق
باللوحة). ومنه $H_{3\times3} = A_8$، ويبيّن استدلال
الأسئلة 6 و8 و18 --- وهو أيضًا لا يتعلق باللوحة --- أن
الوضعيات القابلة للبلوغ هي بالضبط تلك التي $I = +1$: أي نصف
$9!$، أي $181\,440$.

\textbf{23.} نُسمّي الخانات $0, \dots, n-1$ حول \omterm{def:b2:structures:sn}{الدورة}.
وتبادل النقلةُ الفراغ مع أحد جاريه. ونقرأ
القطع بالترتيب الدوري ابتداءً من موضع بُعيد الفراغ: فنحصل على كلمة $w$
تسرد القطع $n - 1$. وتقديم الفراغ خطوة واحدة
يعوّض $(p, w)$ بالزوج $(p + 1, \rho w)$، حيث $p$ خانة
الفراغ و$\rho$ يدير الكلمة دوريًا بمقدار واحد؛ والنقلة
الرجوعية هي المقلوب. فالترتيب \emph{الدوري} للقطع (أي
الكلمة إلى حدّ الدوران) لا متغيّر إذن. وصنف بلوغ
$(p, w)$ هو مدار التطبيق $g \colon (p, w) \mapsto (p+1,
\rho w)$، وهو عنصر رتبته $\operatorname{lcm}(n, n-1) =
n(n-1)$ في جداء الزمرتين الدوريتين (انسحابات
$\Z/n\Z$ ودورانات مواضع الكلمة $n-1$)، والمضاعف المشترك الأصغر
هو $n(n-1)$ لأن $\gcd(n, n-1) = 1$: فلكل صنف
بالضبط $n(n-1)$ وضعية، ولها جميعًا القلادة نفسها.
والأصناف: $n!\,/\,\bigl(n(n-1)\bigr) = (n-2)!$. ومن أجل $n \geq 5$،
$(n-2)! > 2$: فلا متغيّر التماثل (صنفان في أحسن الأحوال)
أعمى عن كل العائق تقريبًا؛ وثراء اللوحة $4
\times 4$
--- حيث التماثل هو العائق \emph{الوحيد} ---
واقعة هندسية أصيلة، لا شكلية.

\textbf{24.} في اللوحتين ترتيب القطع مقلوب كليًا،
ومنه $N = \binom{15}{2} = 105$ في الحالتين (فكل زوج من القطع
منقلب). \emph{الفراغ في موطنه:} $r = 1$، و$N + r = 106$ زوجي:
فهي غير قابلة للحلّ. \emph{الفراغ في الخانة $1$:} الفراغ في السطر
الأعلى، و$r = 4$، و$N + r = 109$ فردي: فهي قابلة للحلّ. لوحتان لا تختلفان
إلا في موضع الثقب تقعان على جانبَي الجدار.

\textbf{25.} \emph{خاصية التشاكل:} تحوّل ``نقلة واحدة =
\omterm{def:b2:structures:sn}{مبادلة} واحدة'' إلى ``نقلة واحدة = قلب إشارة واحد'' (السؤالان
2 و4)، فتجعل $I$ قابلة للحساب نقلةً نقلة. \emph{لاغرانج:} فرض
$6 \mid \abs{\langle s, t\rangle}$ في المبرهنة المساعدة أ وحدّد
حجم الأصناف الجانبية في استبعاد الرتبة $6$ (السؤال 14).
\emph{التوليد بالدورات ذات الطول $3$:} حوّل ``$H$ تحتوي
ما يكفي من الدورات ذات الطول $3$'' إلى ``$H$ تحتوي $A_{15}$ كلها''
(السؤال 16). \emph{المقارنة:} صنعت الدورات المتتالية
الخمس عشرة ذوات الطول $3$ انطلاقًا من جولة $2 \times 2$ واحدة
منقولة بالجولة الكبرى (السؤالان 12--13)، وصنعت
الدورات ذوات الطول $3$، أي $(a\ b\ w)$، في المبرهنة المساعدة ب. \emph{المبدأ الشامل:}
لا متغيّرٌ يبرهن على الاستحالة، وإنشاءٌ صريح يبرهن على
الإمكان، ولا تُحلّ المسألة تمامًا إلا حين يلتقي الحدّان
--- وهنا يلتقيان عند النصف.
\end{solution}
